% ED26-SCI-NASA7009B2024-001: source-specific methodology and explicit inference cases.
\chapter{Simulation Results as Evidence}
\label{sec:p12-simulation-results-as-evidence}

A simulation result concerns a represented system. Using that result to
answer a question about an actual system requires an argument connecting
the representation, its implementation, the supplied inputs and the
observations used to assess it. A proof about the represented system
establishes its conditional mathematical conclusion; the connection to
the target has its own evidence.

This treatment reconstructs the selected methodology of
\emph{NASA-STD-7009B, Standard for Models and Simulations}, approved
5 March 2024, which supersedes revision A with Change 1
\cite[title, Sections 1.1--1.3 and 3.2]{NASA7009B2024}.
Its requirements govern the NASA applications in which the standard is
made applicable; they are not axioms of inference or additional production
requirements for this handbook. The exact decay calculations and
counterexamples below are mathematical illustrations. The separate
BoxBOD data are a published NIST numerical-reference dataset, whose
observational and calibration status is stated explicitly.

\section{Specify the target, model and use}

\begin{hbdefinition}{def:p12-simulation-objects}{The objects in a simulation inference}
The \emph{target system} is the system, phenomenon or process about
which the reader seeks a conclusion. NASA calls it the real-world
system. Its \emph{environment} contains relevant external conditions
and interactions. A \emph{conceptual model} specifies the abstractions,
assumptions, components and relationships chosen to represent that
target. A \emph{mathematical model} supplies the equations, initial and
boundary conditions, and modeling data. A \emph{computational model}
implements the conceptual model through its mathematical, numerical,
logical or qualitative representations. A \emph{simulation} imitates
the selected behavior by using that representation in a specified
\emph{scenario}: inputs, environment, initial conditions and timeline.

A \emph{referent} is the data, information, theory or other representation
against which results are compared. It may be the target, a representative
experiment, a related system or a higher-fidelity model. These are
different sources of evidence. The match between a referent and the
target must itself be justified.
\cite[Section 3.2]{NASA7009B2024}
\index{simulation}\index{conceptual model}\index{computational model}
\index{referent}
\end{hbdefinition}

\begin{hbdefinition}{def:p12-simulation-use}{Five use descriptions and three domains}
In the standard, an \emph{intended use} is the expected purpose and
application; a \emph{permissible use} is one formally allowed;
a \emph{proposed use} is a particular desired application;
an \emph{accepted use} is a proposed use found sufficient by its use
assessment; and \emph{actual use} describes what is or was done.
These labels do not identify the same record.

The \emph{model limits} bound operation according to design and
implementation. The \emph{verification domain} concerns inputs for
which the solution has been assessed against the model's requirements
and accuracy criteria. The \emph{validation domain} concerns inputs
for which relevant outputs compare favorably with the selected
referent. It can differ by output. The proposed-use domain specifies
the conditions needed for the current decision. A program running
within its operational limits does not thereby establish validation
throughout those limits.
\cite[Sections 3.2, 4.2.1.6--4.2.1.7 and E.3.4]{NASA7009B2024}
\end{hbdefinition}

The inference begins with a specific quantity, population or regime,
time horizon and decision tolerance. For example, predicting the sign
of a temperature change, its magnitude within one degree and the
probability of exceeding a threshold require different accuracy and
uncertainty information. A model suitable for one may be unsuitable
for another even if its code and stored coefficients are unchanged.
An abstraction omits or transforms features of the target; its relevant
assumption is that these choices do not materially undermine the
intended use. That assumption is investigated, not supplied by the
act of writing down the model.

\begin{hbmethod}{def:p12-simulation-development-record}{The selected NASA development records}
The following are attributed requirements of the selected standard:
\begin{enumerate}
\item Record units and vector coordinate frames where applicable
      (4.2.1.3, M\&S 46). Conversion and coordinate conventions belong
      to the representation, not just the printed output.
\item Record assumptions and abstractions, their rationales and their
      consequences for intended uses (4.2.1.4, M\&S 11).
\item Record the concepts, structures and mathematics: equations,
      represented behavior or states, and flow of execution
      (4.2.1.5, M\&S 12).
\item Record design/implementation limits, including boundary
      conditions (4.2.1.6, M\&S 13).
\item Record permissible uses in light of intended use, assumptions,
      limits and the verification and validation domains
      (4.2.1.7, M\&S 14).
\end{enumerate}
For off-the-shelf or legacy models, the standard allows relevant
published user guides to supply some of the assumptions and
specification record. A reference supplies only what the referenced
document actually establishes. \cite[Section 4.2.1]{NASA7009B2024}
\end{hbmethod}

\section{Verification, validation and calibration}

\begin{hbdefinition}{def:p12-simulation-vv}{What each comparison establishes}
\emph{Implementation verification} investigates whether the
computational model implements its design and specification.
\emph{Solution verification} investigates errors and achieved
accuracy in a particular computation, including discretization,
iterative convergence, arithmetic and input-handling mistakes.
\emph{Conceptual validation} investigates whether the conceptual
model adequately represents the target for its intended uses.
\emph{Empirical validation} compares an operating model's responses
with appropriate target or referent evidence for those uses.
\emph{Calibration} adjusts model parameters or structure to improve
agreement with a referent. Agreement on data used for that adjustment
is calibration performance; calling it independent predictive
validation requires additional evidence.
\cite[Sections 3.2 and 4.2.3--4.2.6]{NASA7009B2024}
\index{verification!simulation}\index{validation!simulation}
\index{calibration!simulation}
\end{hbdefinition}

The standard requires verification and a record of its domain
(M\&S 15--16), and validation and a record of its domain
(M\&S 17--18). Its recommendations separately call for the methods,
numerical-error estimates, status and unverified aspects; and for
validation methods, experimental design, metrics, referents, datasets,
studies, outcomes and unvalidated aspects. The evidence used as a
validation referent also needs assessment of acquisition, processing,
relevance and accuracy, with the rigor used for the respective model's
validation. The standard permits verification and
validation to occur cyclically or in an interleaved process; it does
not turn one into evidence for the other's distinct question.
\cite[Sections 4.2.3--4.2.6]{NASA7009B2024}

\begin{hbdefinition}{def:p12-simulation-error-chain}{An explicit scalar error chain}
Fix a query condition \(x\). Let \(Y(x)\) be the target quantity,
\(m(x;\theta_*)\) the exact model at target input/parameter values,
\(m(x;\widehat\theta)\) the model at supplied values,
\(m_h(x;\widehat\theta)\) a discretized mathematical calculation,
and \(\widehat m_h\) its computed value. Suppose nonnegative
bounds are justified for the respective differences:
\[
\begin{aligned}
 |Y-m(x;\theta_*)|&\leq E_{\rm model},&
 |m(x;\theta_*)-m(x;\widehat\theta)|&\leq E_{\rm input},\\
 |m(x;\widehat\theta)-m_h(x;\widehat\theta)|&\leq E_{\rm disc},&
 |m_h(x;\widehat\theta)-\widehat m_h|&\leq E_{\rm arith}.
\end{aligned}
\]
These are stated worst-case bounds for the displayed query, not
automatically confidence intervals or estimates of probabilities.
An unavailable bound remains unavailable rather than being set to
zero. A different decomposition is possible, but its terms must
refer to explicit differences and must not double-count them.
\end{hbdefinition}

\begin{hbtheorem}{thm:p12-simulation-error-chain}{proved=annex:p12.simulation.error-chain}{A conditional numerical-to-target certificate}
Under the four displayed premises,
\[
 |Y(x)-\widehat m_h|\leq
 E_{\rm model}+E_{\rm input}+E_{\rm disc}+E_{\rm arith}.
\]
This bound can be used in the sign/order certificates of
\cref{thm:p12-numerical-certification}. Reducing only discretization
or arithmetic error leaves the other terms in the conclusion.
The inequality proves no premise about target correspondence.
\end{hbtheorem}

\begin{hbtheorem}{thm:p12-simulation-reference-comparison}{proved=annex:p12.simulation.reference-comparison}{What a comparison with uncertain reference data implies}
Fix nonnegative bounds \(E_r,E_n\) and tolerance \(\tau\).
At one query condition let a recorded referent value \(r\) satisfy
\(|r-Y|\leq E_r\), and let the computed model value satisfy
\(|\widehat m-m|\leq E_n\). Then
\[
 \max\{0,|\widehat m-r|-E_n-E_r\}
 \leq |m-Y|\leq |\widehat m-r|+E_n+E_r.
\]
Consequently \(|\widehat m-r|+E_n+E_r\leq\tau\) suffices to
establish \(|m-Y|\leq\tau\) at this point. Conversely,
\(|\widehat m-r|-E_n-E_r>\tau\) disproves that tolerance.
The hypotheses must concern the same quantity, units, condition and
referent-to-target relation.
\end{hbtheorem}

NASA's Appendix E.3.4 discusses accuracy and uncertainty together and
requires at least some interval overlap for its favorable-comparison
assessment. Overlap alone need not meet a decision tolerance. Two
estimates both reported as \(0\pm10\) overlap, while target and
model values \(-10\) and \(10\) are both compatible and differ by
twenty. Enlarging an uncertainty interval can create overlap without
creating useful predictive accuracy. The comparison criterion must
be fixed for the intended use, with its interpretation stated.

\section{A decay calculation with an exact mathematical reference}

\begin{hbdefinition}{def:p12-simulation-decay}{The model and its Euler calculation}
Let \(q_0\geq0\), \(\lambda\geq0\), and consider
\[
 q'(t)=-\lambda q(t),\qquad q(0)=q_0,\quad t\geq0.
\]
Time has a declared unit, \(\lambda\) its reciprocal unit, and
\(q\) the quantity's unit. The product \(\lambda t\) is
dimensionless. This is a specified mathematical model; application
to a physical decay process needs evidence for its constant-rate and
other modeling assumptions. For a step \(h>0\), define exact-real
Euler iterates \(v_0=q_0\),
\(v_{j+1}=(1-\lambda h)v_j\) for integers \(j\geq0\).
The step count \(N\) is a nonnegative integer; zeroth powers in the
closed form denote one, including a zero multiplier. These iterates are a discrete
mathematical algorithm, before any machine rounding.
\end{hbdefinition}

\begin{hbtheorem}{thm:p12-simulation-decay-euler}{proved=annex:p12.simulation.decay-euler}{Exact solution, qualitative behavior and step error}
The differential equation has the unique differentiable solution
\(q(t)=q_0\exp(-\lambda t)\), and
\(v_N=q_0(1-\lambda h)^N\). If \(q_0>0\), the iterates are
nonnegative and nonincreasing when \(0\leq\lambda h\leq1\).
For \(1<\lambda h<2\) they alternate in sign with decreasing
magnitude; at \(\lambda h=2\) their magnitude stays constant,
and for \(\lambda h>2\) it grows.
If \(0<\lambda h<1\) and \(T=Nh\), then
\[
 0\leq q(T)-v_N\leq
 q(T)\frac{\lambda^2 T h}{2(1-\lambda h)}.
\]
For \(\lambda=0\), the solution and iterates are constant exactly.
\end{hbtheorem}

\begin{hbexample}{def:p12-simulation-step-table}{A verified approximation can still answer the wrong physical question}
For the dimensionless example \(q_0=\lambda=T=1\),
\(q(1)=\exp(-1)\approx0.36787944\). The following values use
exact-real iterates; the decimals are rounded for display.
\begin{center}
\begin{tabular}{rrrr}
\toprule
Steps \(N\)&\(v_N\)&Actual discretization error&Proved upper bound\\
\midrule
4 &0.31640625&0.05147319&0.06131324\\
10&0.34867844&0.01920100&0.02043775\\
20&0.35848592&0.00939352&0.00968104\\
\bottomrule
\end{tabular}
\end{center}
The bound at twenty steps proves an absolute error below \(0.01\)
for this differential equation. It supplies no measurement of a
physical system. For \(h=2,\lambda=1\), even exact arithmetic
returns \(v_1=-q_0\): the numerical method violates the model's
positive-decay behavior. That failure is present before roundoff.
\end{hbexample}

\begin{hbtheorem}{thm:p12-simulation-arithmetic-propagation}{proved=annex:p12.simulation.arithmetic-propagation}{Propagation of separately bounded arithmetic errors}
Let \(a=\lambda h\in[0,1]\), and suppose a computed recurrence is
\[
 \widehat v_{j+1}=(1-a)\widehat v_j+\delta_j,
 \quad |\delta_j|\leq\rho_j,\quad
 |\widehat v_0-q_0|\leq E_0.
\]
Then
\[
 |\widehat v_N-v_N|\leq
 (1-a)^N E_0+\sum_{j=0}^{N-1}(1-a)^{N-1-j}\rho_j.
\]
If every \(\rho_j\leq\rho\), this is at most
\((1-a)^N E_0+\rho[1-(1-a)^N]/a\) for \(a>0\),
and at most \(E_0+N\rho\) for \(a=0\).
Coefficient rounding and every executed operation must be included
in the actual \(\delta_j\) budget. The recurrence does not assume
that these errors are random or independent.
\end{hbtheorem}

\begin{hbtheorem}{thm:p12-simulation-decay-inputs}{proved=annex:p12.simulation.decay-inputs}{A simultaneous input-sensitivity bound for this model}
For \(q_0,\widetilde q_0\in[0,Q]\),
\(\lambda,\widetilde\lambda\in[0,\Lambda]\), and
\(t,\widetilde t\in[0,T]\),
\[
 |q_0e^{-\lambda t}-\widetilde q_0e^{-\widetilde\lambda\widetilde t}|
 \leq |q_0-\widetilde q_0|
 +QT|\lambda-\widetilde\lambda|
 +Q\Lambda|t-\widetilde t|.
\]
This joint bound follows from this particular model. It does not
justify ignoring interactions in a different sensitivity problem.
\end{hbtheorem}

\section{Published observations and a fitted reference curve}

\begin{hbdefinition}{def:p12-simulation-boxbod}{The NIST BoxBOD evidence record}
The NIST Statistical Reference Datasets collection supplies six
BoxBOD observations, identifying \(t\) as incubation time in days
and \(y\) as biochemical oxygen demand in mg/l. Its reference
model is \(y=B(1-\exp(-\lambda t))+e\), and it reports the
parameter estimates
\[
 \widehat B=213.80940889,\qquad
 \widehat\lambda=0.54723748542.
\]
The first parameter has units mg/l and the second reciprocal days.
The data file attributes the observations to Box, Hunter and Hunter
(1978), pp.~483--487. Here the inspected primary record is NIST's
complete ASCII dataset; no additional experimental protocol,
measurement-error bound or independent validation sample is supplied
by that record. \cite{NISTBoxBOD}

NIST's certification methodology says the reference calculations
used 128-bit precision and were confirmed with at least two algorithms
and software packages using analytic derivatives. This concerns
the numerical reference solution to the fitting problem. It does
not certify that the fitted curve is a physical law or predictively
adequate for every use. \cite[Methodology]{NISTBoxBODCertification}
\end{hbdefinition}

Using the reported parameters in the stated expression gives the
following approximate predictions and residuals. No new fit or new
experiment is being reported.
\begin{center}
\begin{tabular}{rrrr}
\toprule
Time (days)&Recorded BOD&Fitted value&Recorded minus fitted\\
\midrule
1 &109&90.1109 &18.8891\\
2 &149&142.2441&6.7559\\
3 &149&172.4056&-23.4056\\
5 &191&199.9509&-8.9509\\
7 &213&209.1708&3.8292\\
10&224&212.9111&11.0889\\
\bottomrule
\end{tabular}
\end{center}
All response columns use mg/l. The residual sum of squares is
approximately \(1168.0088766\), matching NIST's reported reference
value at the displayed precision. NIST reports four residual degrees
of freedom and residual standard deviation approximately \(17.0880724\).
That fitted residual statistic is not a supplied deterministic
measurement-error bound. The six observations were used for the
reference fit; reusing them does not create an independent assessment
of predictive transfer. The evidence distinctions in
\cref{def:p19-evidence-record} apply here as well.

\begin{hbtheorem}{thm:p12-simulation-boxbod-residual}{proved=annex:p12.simulation.boxbod-residual}{A numerical-reference success need not meet a response tolerance}
With the two reported decimal parameters treated as exact for this
calculation, the fitted value at day three exceeds the recorded
value \(149\) by more than \(22.04\) mg/l. Therefore an illustrative
criterion requiring discrepancy at every recorded point to be at
most \(10\) mg/l fails. Ten mg/l is a declared comparison criterion
for this example, not a NASA or NIST acceptance threshold.
This failure against a recorded value does not determine the
unreported measurement uncertainty or prove which physical mechanism
generated that observation.
\end{hbtheorem}

The fitted curve's remaining deficit
\(q(t)=\widehat B-\widehat B(1-\exp(-\widehat\lambda t))\)
obeys the exact decay model above. Thus one can verify a numerical
evaluation against its analytic expression, characterize its step
and arithmetic errors, and still face a separate question about
agreement with observations. A third use, predicting an unobserved
system or a changed operating regime, needs additional evidence.

\section{Transfer and the limits of agreement}

\begin{hbtheorem}{thm:p12-simulation-finite-comparisons}{proved=annex:p12.simulation.finite-comparisons}{Finite comparisons alone give no uniform prediction bound}
For \(n\geq1\) distinct real comparison inputs \(t_1,\ldots,t_n\),
a real point \(s\) outside that finite set, and any continuous model
function \(f:\mathbb R\to\mathbb R\), there are continuous functions agreeing with \(f\)
at all the comparison inputs but differing from it at \(s\) by
an arbitrarily large amount. If \(f\) is smooth or analytic, the
alternative can retain that property. This remains true when
\(s\) is between comparison inputs.
\end{hbtheorem}

The standard calls the region enclosing favorable comparisons a
validation domain. Naming that region does not, as a mathematical
inference, supply a uniform error bound at every interior point.
Interpolation needs an additional structural or statistical argument;
extrapolation also needs a warrant for the new regime. Output-specific
validation matters: agreement for one output does not establish
agreement for every other quantity derived from the simulation.
\cite[Sections 3.2 and E.3.4]{NASA7009B2024}

\begin{hbtheorem}{thm:p12-simulation-regime-change}{proved=annex:p12.simulation.regime-change}{Exact agreement before a change does not identify later decay}
Fix \(T>0\), \(q_0>0\), and \(\lambda,\mu>0\) with
\(\lambda\ne\mu\). The two functions
\[
 q_1(t)=q_0e^{-\lambda t},\qquad
 q_2(t)=
 \begin{cases}
 q_0e^{-\lambda t},&0\leq t\leq T,\\
 q_0e^{-\lambda T}e^{-\mu(t-T)},&t>T
 \end{cases}
\]
are continuous, positive and nonincreasing, and agree throughout
\([0,T]\). They differ at every \(t>T\). Therefore even complete,
exact knowledge of the earlier trajectory cannot distinguish them
without a premise about persistence of the rate mechanism.
\end{hbtheorem}

For \(q_0=\lambda=T=1,\mu=2\), the values at time two are
\(e^{-2}\approx0.135335\) and \(e^{-3}\approx0.049787\).
They lie on opposite sides of \(0.1\), despite identical earlier
trajectories. This is a constructed counterexample, not a report of
a physical experiment. Smaller time steps in the first model
converge more accurately to its own wrong-regime prediction if the
second trajectory is the target.

\section{Uncertainty, pedigree and the separate credibility factors}

\begin{hbdefinition}{def:p12-simulation-uncertainty}{The source's uncertainty distinctions}
NASA distinguishes \emph{aleatory uncertainty}, variation inherent
in the physical system that is not removed without changing that
system or its operation, from \emph{epistemic uncertainty}, missing
knowledge about quantities or processes. Its \emph{model uncertainty}
concerns variation caused by assumptions, formulas and representations;
\emph{numerical errors} concern finite precision and numerical
approximations. These distinctions do not by themselves assign a
probability distribution or show that a chosen propagation method
covers every relevant source.

\emph{Data pedigree} traces data used in model development from
source through transmission, storage and processing. \emph{Input
pedigree} traces the data supplied for a particular use. Calibration
datasets, runtime inputs and validation referents therefore need
their own identities and transformation records. A metadata record
describes an artifact; it is not the underlying measurement or model.
\cite[Section 3.2 and Sections E.3.2, E.4.3--E.4.4]{NASA7009B2024}
\end{hbdefinition}

During development, requirements M\&S 19 and 21 call for the processes
and rationales used to characterize uncertainty in referent data, and
the qualitative or quantitative uncertainty incorporated into the model.
The recommendation in 4.2.8 records significant processes, effects,
scenarios or environments left out. During use, M\&S 28--29 extend
the record to inputs, results and quantities derived from results,
including the treatment of model-uncertainty factors and factors of
safety; 4.3.4 again recommends recording significant omissions.
The standard's semi-quantitative model-uncertainty factor is an
experience-based adjustment. Its magnitude is not automatically
a proved error bound, and a margin or safety factor needs a stated
reason to cover the uncertainties assigned to it.
\cite[Sections 4.2.7--4.2.8 and 4.3.3--4.3.4]{NASA7009B2024}

Revision B separates an assessment of model \emph{capability} at
development/release from an assessment of the \emph{results} of a
particular use. Both are reported. NASA uses \emph{credibility} for
the quality that elicits belief or trust in the results; its factor
assessment is evidence for a use-dependent judgment, not a probability
assigned to the truth of a simulation.
The five capability factors are data pedigree, verification,
validation, development technical review, and development
process/product management. The six results factors are use
assessment, input pedigree, uncertainty characterization, results
robustness, use/analysis technical review, and use process/product
management. The selected evidence factors
have the following full-level distinctions; zero denotes insufficient
evidence, not a numerical probability of model failure.
\cite[Section 3.2 and Appendix E.1--E.4]{NASA7009B2024}

\begin{description}
\item[Development data and runtime input pedigree (Tables 7 and 16).]
Level 1 has some informal traceability and known sources for all
significant data; uncertainty can remain unestimated. Level 2 has
most data known and formally traceable, known processes for significant
data, and at least estimated uncertainties for all data. Level 3
requires all data traced to a sufficiently representative referent,
with all significant data acceptable in accuracy, precision and
uncertainty. Level 4 requires all data traced to the real-world
system and acceptable in those qualities. The two factors assess
different datasets even though their level structures correspond.
\item[Verification (Table 8).]
Level 1 is informal verification with some documentation or evidence
of completeness/success. Level 2 uses documented practices across
components, features, capabilities and couplings, with most important
model errors meeting requirements. Level 3 uses rigorous practices
on the end-to-end model, with all important errors meeting the
requirements. Level 4 uses reliable practices and error-estimation
methods, with all model errors meeting the specified requirements.
These are evidence assessments, not a supplied universal program proof.
\item[Validation (Table 9).]
Level 1 requires conceptual validation. Level 2 adds favorable
comparison with a representative system, or a higher-fidelity model
itself meeting level 2, at validation points in the operating domain
for at least some important outputs. Level 3 uses the actual system
in a representative environment, or an appropriately assessed
higher-fidelity model; the points significantly span the operating
domain and all important outputs compare favorably. Level 4 uses
the actual operating environment, or a higher-fidelity model meeting
level 4; points completely span the domain and all response
quantities compare favorably. A level above 1 retains its conceptual
validation prerequisite. Agreement with another model requires
that model's own target-correspondence evidence.
\item[Uncertainty characterization (Table 17).]
Level 1 identifies input sources qualitatively, without addressing
their propagation to outputs. Level 2 identifies most sources,
quantifies and classifies them, and uses reduced-order or
reduced-dimension propagation. Level 3 provides quantitative
uncertainty for most outputs after propagation of all known sources.
Level 4 adds statistical analysis for all outputs after rigorous,
validated propagation of all known sources. The source lists ANOVA
as an example, not as an appropriate universal analysis for every
simulation. Unknown omitted sources remain a separate issue.
\item[Results robustness (Table 18).]
Level 1 estimates sensitivity by analogy with quantified similar
problems. Level 2 quantifies a few sensitivities, identifies few or
none of the most sensitive variables and has no known combination
sensitivities. Level 3 covers many variables, many key sensitivities
and some combinations. Level 4 covers most variables, most or all
key sensitivities, and many combinations. The source's percentage
descriptors for few/many/most are guidelines rather than universal
cutoffs. This factor measures how well sensitivities are known;
a well-characterized system can still be highly sensitive.
\end{description}

Technical-review and process/product-management factors remain part
of NASA's complete capability/results assessments. This chapter's
selected comparison does not replace those assessments with just
the factors above. Where a level lists several conditions, the
source requires them all unless it explicitly gives alternatives.
It specifies discrete attainment rather than fractional credit.
The illustrative 0--4 labels are ordinal. Appendix E.2 discourages
comparison or combination of levels across different factors;
each factor can still be compared with its own established threshold.
The labels do not constitute a probability scale or an overall truth score.

\begin{hbtheorem}{thm:p12-simulation-ordinal-scores}{proved=annex:p12.simulation.ordinal-scores}{An ordinal label does not determine an arithmetic aggregate}
A ranking of factor profiles obtained by averaging their numerical
labels need not survive a strictly increasing relabeling of the
same levels. Therefore the factor order alone does not identify
such an aggregate. An additional justified aggregation or decision
model would supply information absent from the ordinal labels.
\end{hbtheorem}

Sensitivity analysis in 4.3.5 investigates responses to input changes
and records its extent and results (M\&S 30). Validation must also
investigate whether the modeled response sensitivities correspond
to those of the target. A numerical derivative of a wrong model
can be computed accurately. Local invariance at one operating point
does not establish invariance across a larger parameter region.
\cite[Section 4.3.5 and Appendix E.4.5]{NASA7009B2024}

\section{A use assessment and an accountable conclusion}

\begin{hbmethod}{def:p12-simulation-use-record}{The selected NASA use and reporting obligations}
For the applications to which the standard applies, the selected
clauses distinguish requirements from recommendations:
\begin{enumerate}
\item Requirements 4.3.1.1--4.3.1.4 (M\&S 22--25) record the
      proposed use, the assessment against permissible uses, input
      pedigree, and setup/execution processes and rationales.
\item Requirement 4.3.1.5 (M\&S 26) calls either for use within
      permissible uses and verification/validation domains, or for
      explicit warnings identifying the exceeded limit, its extent
      and assessed consequences. A warning reports an unsupported
      extension; it does not validate the extension. M\&S 27
      records observed warnings, errors and their explanations.
\item Recommendations 4.3.2 record relevant target characteristics,
      model revisions, calibrations and calibration domain, input
      preparation data/software, analysis procedures, result versions,
      relevant prior uses, and failure modes and messages.
\item Requirements 4.3.6.1--4.3.6.2 (M\&S 31 and 49) record the
      six results factors and risks of using the analysis. The latter
      considers criticality, review findings, defects, caveats,
      uncertainty and its methods, capability/results assessments,
      qualifications, and significant changes from prior releases
      or uses. Section 4.3.7 recommends recording the justification
      of each assessed factor.
\item Requirement 4.3.8.1 (M\&S 32) reports unmet acceptance
      criteria, violated assumptions or model limits, execution
      warnings/errors, unfavorable use assessments, setup/use
      issues, waivers and outstanding defects, with at least a
      qualitative account of their impacts. Missing supporting
      evidence for the referenced requirements also calls for a warning.
\item Requirements 4.3.8.2--4.3.8.3 (M\&S 33--34) report a
      quantitative uncertainty estimate, a justified qualitative
      description, or explicitly that neither is available; they
      also report the processes used to obtain the characterization.
      The explicit-unavailability option supplies no numerical bound.
\item Requirements 4.3.8.4--4.3.8.5 (M\&S 50 and 35) report both
      capability and results assessments, including gaps from
      the established thresholds. Requirements 4.3.8.6--4.3.8.9
      (M\&S 36--39) report review findings, relevant personnel
      qualifications, development/use records, and the risk
      assessment with the rationale for accepting those risks.
\item Recommendations 4.3.9 state whether credibility is sufficient
      for actual use, provide access to supporting records,
      preserve controlled result versions, report departures from
      recommended practice and dissenting technical opinions, and
      communicate serious concerns.
\end{enumerate}
The standard identifies required statements with \emph{shall};
\emph{should} marks a recommendation. Its sufficiency thresholds
are determined for the program/project rather than prescribed as
one universal level by 4.3.6.
\cite[Sections 1.2.4--1.2.5 and 4.3.1--4.3.9]{NASA7009B2024}
\end{hbmethod}

A usable reference conclusion identifies the target question and
actual scenario; the exact model, data and execution versions; which
implementation and solution claims were checked; the relevant
referent and validation comparisons; numerical and empirical
uncertainty, including omissions; sensitivity and transfer premises;
and the decision criterion applied. It then states what follows
conditionally and which parts remain unresolved. Repeating a
deterministic run, adding digits or satisfying a reporting clause
does not supply missing target evidence.
\index{simulation!credibility}\index{simulation!domain of validation}
\index{simulation!uncertainty}\index{simulation!transfer}
